% Storyline 2, paragraph 5: findings and implication.
We find that private observers can agree with one another and still distinguish cooperators from defectors, although they share no information, provided that inertia is allocated in a norm-specific direction. Agreement alone is not sufficient. With the same small inertia in both directions, observers under Shunning agree only because they regard nearly everyone as bad. Compact architectures, selected solely by these two reputational criteria, make cooperation stable and easy to reach under each of four classical norms (Scoring, Shunning, Stern Judging and Simple Standing). At the same depth, inertia placed on the wrong path can leave no cooperation at all, so how evidence is processed matters more than how much of it is retained. Finally, the tolerance that makes private verdicts robust can also shelter opportunists who defect only occasionally, a vulnerability that the aggregation and score-based rules we test share. Allowing warning states to carry partial consequences closes this loophole for the opportunistic strategies we consider, while cooperation remains high. Consensus about reputations therefore need not be imposed through shared information. It can emerge from how each individual weighs its own evidence.
